
\paragraph{Income} The main source is the INE's Household Income Distribution Atlas (\textit{Atlas de Distribución de Renta de los Hogares}, ADRH; \citealt{ine2025adrh}). The ADRH links personal income-tax records from the national tax agency (AEAT) and the foral tax authorities of the Basque Country and Navarre to the population register (\textit{Fichero Precensal de Población}) and attributes income to the census tract where each recipient lives. Income refers to a calendar year and population to 1 January of the following year; people living in collective establishments, such as care homes or barracks, are excluded. I use the latest edition, which refers to income in \BaseYear{}.

Net income is gross income minus the personal income tax due, employees' social security contributions and the wealth tax. Gross income comprises wages, including remuneration in kind and employers' contributions to pension plans; pensions; unemployment benefits; other public benefits, including tax-exempt ones such as disability, dependency and family benefits and regional minimum-income schemes; interest, dividends and other capital income; rental income; business income; and the notional income that the income tax attributes to second homes. Non-filers are covered through the information returns of those who pay them. The ADRH excludes income that the tax authorities do not record, such as informal earnings, transfers between households and foreign pensions, as well as capital gains. It also covers only people living in households with some recorded income, which in 2016 included more than 98.6\% of the population of the common fiscal territory.

Household net income is divided by the household's consumption units on the modified OECD scale (1 for the first adult, 0.5 for every other member aged 14 or over and 0.3 for every member under 14) and the result is assigned to every member. The ADRH publishes thirteen statistics of this distribution of persons: its mean, median, Gini coefficient and P80/P20 ratio, and the shares of the population with income below three fixed amounts (\texteuro 5,000, \texteuro 7,500 and \texteuro 10,000) and below or above six fractions of the national median (40, 50, 60, 140, 160 and 200\%). All the distributions in this note are distributions of persons.

\paragraph{Publication rules} Four features of the ADRH matter for what follows. First, the mean and median of income per consumption unit, the Gini coefficient and the P80/P20 ratio are published only for tracts and municipalities with at least 100 residents, and the shares only for those with at least 500 (\ShareTracts{} tracts, with \ShareTractsPopPct{}\% of the population). Below that size only income per person and per household is published, taken from the district or municipality that contains the tract or, for municipalities with fewer than 100 residents, averaged over the municipalities of that size in the same agricultural district (\textit{comarca agraria}). Second, the Gini coefficient is computed after replacing the lowest and highest 1.5\% of incomes with the nearest remaining value. Third, published tract indicators are bounded at the population-weighted 0.1st and 99.5th percentiles across tracts; in \BaseYear{} this caps the median of \CapTracts{} tracts, home to \CapPopPct{}\% of the population, at \texteuro\CapMedian{}, and the Gini coefficient of \CapGiniTracts{} tracts at \CapGini{}. Fourth, outside Navarra, published medians are multiples of \texteuro 350 (\MedBinnedPct{}\% of tracts), consistent with the midpoints of \texteuro 700 intervals, and the P80/P20 ratio is rounded to one decimal.

\paragraph{Other sources} To project incomes beyond \BaseYear{} (Section \ref{sec:nowcast}), I use the annual tax revenue report of the Spanish tax agency (\textit{Informe Anual de Recaudación Tributaria}, Agencia Tributaria). Its Table 2.1 reports household income from tax sources (withholdings and personal income-tax returns, the records the ADRH is built from) by type of income, together with the income tax accrued, about six months after the end of the year; the latest two years are provisional. Population is the national accounts series of Eurostat (\texttt{nama\_10\_pe}). The municipal indicators that \href{https://comparatuingreso.es/}{\mbox{comparatuingreso.es}} shows alongside its results (the share of the population aged 15 and over with higher education and the share born abroad) come from the INE's annual population census and play no role in the estimation.

\paragraph{Territorial units} Spain is divided into 17 autonomous communities, two autonomous cities (Ceuta and Melilla), 52 provinces and, in the ADRH, \MunAdrh{} municipalities. For statistical purposes, municipalities are divided into districts and districts into census tracts (\textit{secciones censales}), which typically have between 1,000 and 2,500 residents.

\paragraph{Summary statistics} Table \ref{tab:summary} describes the \TractsTotal{} tracts in \BaseYear{}. The average tract has \TractPopMean{} residents (standard deviation \TractPopSD{}). Mean net income per consumption unit averages \texteuro\TractEquivMean{} across tracts, \EquivOverPcPct{}\% more than income per person (\texteuro\TractPcMean{}), and the Gini coefficient averages \TractGiniMean{} points. The average tract has a dependency ratio of \TractDepRatioMean{}, a mean age of \TractMeanAge{} years and \TractSingleHhMean{}\% single-person households.

\begin{table}[H]
\centering
\begin{threeparttable}
\onehalfspacing
\caption{Summary statistics, census tracts}\label{tab:summary}
\footnotesize
\begin{tabular}{@{}lccccc@{}}
\toprule
 & Min & Max & Mean & SD & \% missing \\
\midrule
\multicolumn{6}{@{}l}{\textit{Income distribution}} \\
\quad Net income per person (\texteuro) & 5,522 & 36,918 & 15,025 & 4,432 & 1.8 \\
\quad Net income per consumption unit (\texteuro) & 8,797 & 58,752 & 22,036 & 6,915 & 5.5 \\
\quad Median income per consumption unit (\texteuro) & 7,350 & 44,450 & 19,877 & 5,598 & 5.5 \\
\quad Gini index & 19.7 & 44.2 & 28.6 & 3.9 & 5.5 \\
\multicolumn{6}{@{}l}{\textit{Demographics}} \\
\quad Population & 2 & 12,327 & 1,315 & 678 & 1.6 \\
\quad Dependency ratio & 0.00 & 4.00 & 0.63 & 0.18 & 1.6 \\
\quad Mean age & 27.9 & 74.8 & 45.6 & 5.5 & 1.6 \\
\quad Single-person households (\%) & 0.0 & 100.0 & 30.7 & 9.6 & 1.6 \\
\midrule
Census tracts & \multicolumn{5}{l}{37,072} \\
\bottomrule
\end{tabular}

\begin{tablenotes}[para,flushleft]
\scriptsize
\textbf{Notes}: Unweighted statistics across census tracts. Income in \BaseYear{} euros. The dependency ratio is the population under 18 and over 65 relative to the population aged 18 to 65. Source: INE, ADRH, and author's calculations.
\end{tablenotes}
\end{threeparttable}
\end{table}

\paragraph{Coverage} \GiniMissTracts{} tracts (\GiniMissPct{}\%) lack the mean and median per consumption unit, the Gini coefficient and the P80/P20 ratio. They are tracts below the 100-resident threshold: the median one has \GiniMissMedianPop{} residents, whereas every tract with a published Gini coefficient has at least \ObsGiniMinPop{}. Together they house \GiniMissPopPct{}\% of the population. For \GiniImputedTracts{} of them the ADRH publishes income per person. I impute their mean income per consumption unit as income per person times the population-weighted ratio of the two measures in the province, and their Gini coefficient as described in Section \ref{sec:gini}. The remaining \TractsDropped{} tracts have no income data and are dropped. After imputation, the share of tracts without income per consumption unit falls from \MissEquivPct{}\% to \MissEquivAfterPct{}\%.
